consequence, the immersions @@M0@@ (@@M1@@) are regular and remain regular immersions @@M2@@ under the base change @@M3@@. Note that, by the exact homotopy sequence (IX 1.1), @@M4@@ comes from an element @@M5@@, so the element @@M6@@ of @@M7@@ comes from the element @@M8@@ of @@M9@@, and consequently its image in @@M10@@ under (7.3.1) is precisely
\sgaStableTarget{sga6.e10.eq.unnamed.occ2}@@M11@@
Of course, for a section @@M12@@ of @@M13@@ over @@M14@@, one puts
@@M15@@
is defined as above, and can be regarded as obtained from @@M16@@ by the flat base change @@M17@@; it is therefore a regular immersion, a fortiori a perfect morphism, which gives a meaning to the expression @@M18@@. For the same reason, the expression @@M19@@ is defined, and, by transitivity of inverse images (IV 2.12) in the following commutative diagram of perfect morphisms:
@@M20@@
one finds that the left-hand side of @@M21@@ is equal to
\sgaStableTarget{sga6.e10.eq.unnamed.occ3}@@M22@@
Since @@M23@@ is algebraically closed, the reduced irreducible components of @@M24@@ are geometrically integral over @@M25@@ and their intersection points are rational over @@M26@@. One concludes readily that the expression @@M27@@ is algebraically equivalent to zero in @@M28@@, using the connectedness
